\section{Phản ứng cháy của hydrocarbon}

\subsection{Cháy hoàn toàn}

Khi hydrocarbon cháy hoàn toàn trong oxygen dư, sản phẩm là \ce{CO2} và \ce{H2O}.

\begin{ghinho}{Một số phương trình cần viết thành thạo}
\begin{tasks}(2)
\task $\ce{CH4 + 2O2 -> CO2 + 2H2O}$.
\task $\ce{2C2H6 + 7O2 -> 4CO2 + 6H2O}$.
\task $\ce{C3H8 + 5O2 -> 3CO2 + 4H2O}$.
\task $\ce{2C4H10 + 13O2 -> 8CO2 + 10H2O}$.
\end{tasks}
\end{ghinho}

Với alkane \ce{C_nH_{2n+2}}:
\[
\ce{C_nH_{2n+2} + \frac{3n+1}{2}O2 -> nCO2 + (n+1)H2O}.
\]

\subsection{Đọc thông tin từ sản phẩm cháy}

Nếu một hydrocarbon cháy hoàn toàn thì toàn bộ carbon của hydrocarbon đi vào \ce{CO2}, còn hydrogen đi vào \ce{H2O}. Vì vậy:
\[
n_C=n_{\ce{CO2}},\qquad n_H=2n_{\ce{H2O}}.
\]

Đây là cơ sở để xác định thành phần hydrocarbon từ sản phẩm cháy.

\subsection{Tính theo phương trình cháy}

Vẫn sử dụng bốn bước quen thuộc:
\begin{enumerate}
\item Tính số mol.
\item Viết phương trình hoặc sơ đồ phản ứng.
\item Tính các giá trị trung gian cần thiết.
\item Tính đại lượng đề hỏi và kết luận.
\end{enumerate}

\begin{vidu}
Đốt cháy hoàn toàn \SI{4.48}{\liter} propane ở điều kiện có $V_m=\SI{24.8}{\liter\per\mole}$. Tính thể tích oxygen cần dùng và thể tích \ce{CO2} tạo thành ở cùng điều kiện.
\end{vidu}
\begin{loigiai}
\textbf{Bước 1. Tính số mol}
\[
n_{\ce{C3H8}}=\frac{4.48}{24.8}=\SI{0.1806}{\mole}.
\]

\textbf{Bước 2. Viết phương trình hóa học}
\[
\begin{array}{c@{\;}c@{\;}c@{\;}c@{\;}c}
\ce{C3H8} & + & \ce{5O2} & \ce{->} & \ce{3CO2 + 4H2O}\\
0.1806 & & 0.9030 & & 0.5418
\end{array}
\]

\textbf{Bước 3. Tính các giá trị trung gian}
\[
n_{\ce{O2}}=\SI{0.9030}{\mole},\qquad n_{\ce{CO2}}=\SI{0.5418}{\mole}.
\]

\textbf{Bước 4. Tính đại lượng đề hỏi}
\[
V_{\ce{O2}}=0.9030\times24.8\approx\SI{22.39}{\liter},
\]
\[
V_{\ce{CO2}}=0.5418\times24.8\approx\SI{13.44}{\liter}.
\]

\textbf{Vậy,} cần khoảng \SI{22.39}{\liter} \ce{O2} và tạo khoảng \SI{13.44}{\liter} \ce{CO2}.
\end{loigiai}

\begin{tuluyen}
Viết phương trình cháy hoàn toàn của các chất sau:
\begin{tasks}(2)
\task methane;
\task ethane;
\task propane;
\task butane;
\task ethylene;
\task một alkane có công thức \ce{C5H12}.
\end{tasks}
\dapso{Cân bằng theo bảo toàn C, H rồi O.}
\end{tuluyen}

\begin{tuluyen}
Đốt cháy hoàn toàn \SI{0.25}{\mole} butane. Tính số mol \ce{O2} cần dùng, số mol \ce{CO2} và \ce{H2O} tạo thành.
\dapso{$n_{\ce{O2}}=\SI{1.625}{\mole}$; $n_{\ce{CO2}}=\SI{1.00}{\mole}$; $n_{\ce{H2O}}=\SI{1.25}{\mole}$.}
\end{tuluyen}
